\section{Discussão}

Este capítulo interpretará os resultados em relação à pergunta de pesquisa, aos objetivos, às hipóteses de trabalho e à literatura. A discussão deverá manter separadas a evidência técnica produzida, suas implicações para o caso estudado e as condições institucionais necessárias para eventual uso na administração pública.

\subsection{Recuperação, geração e verificabilidade}

% TODO: discutir a hipótese de que recuperação pertinente é necessária, mas não
% suficiente, identificando mecanismos e casos que sustentem ou contrariem isso.

\subsection{Desempenho por categoria de pergunta}

% TODO: discutir diferenças entre perguntas simples, multietapas, comparativas,
% numéricas e não respondíveis.

\subsection{Limites do paradigma LLM-as-a-judge}

% TODO: interpretar concordâncias e divergências frente à avaliação humana.

\subsection{Implicações para a administração pública}

% TODO: distinguir uso assistivo, prontidão operacional e requisitos de
% governança, curadoria, auditoria, supervisão e contestação.

\subsection{Ameaças à validade e limites de generalização}

% TODO: validade interna, de construto e externa; dependências de corpus,
% configuração, modelos, rubrica, amostra e avaliadores.
